\documentclass[11pt]{article}

\usepackage[margin=1in]{geometry}
\usepackage{amsmath,amssymb,booktabs}
\usepackage[T1]{fontenc}
\usepackage{hyperref}
\usepackage{listings}

\title{DCLAR: Data Center Leak Assessment Robot\\High-Level System Design}
\author{Student Design Prototype}
\date{\today}

\begin{document}
\maketitle

\begin{abstract}
DCLAR is a software prototype for planning a leak-search route over a gridded
data-center floor. It separates floor-tile routing from within-tile coverage,
represents turns as in-place pivots, and stores destinations in a global map
frame. Spill locations can be annotated at sub-tile resolution. The current
program visualizes this plan; it does not read a leak sensor or command physical
motors.
\end{abstract}

\section{System boundary}

The modeled workspace is a $10\times10$ array of $1\,\mathrm{ft}$ square
cells. Rack rows are represented by blocked cells. A simulated robot has a
$7\,\mathrm{in}\times5\,\mathrm{in}$ rectangular footprint; its forward arrow
lies along the $5\,\mathrm{in}$ side, leaving the $7\,\mathrm{in}$ axis
perpendicular to forward motion. The current software uses these dimensions to
space coverage passes, but it does not yet perform continuous collision checks
against rack geometry.

\begin{center}
\begin{tabular}{ll}
\toprule
Subsystem & Responsibility \\
\midrule
Map model & Cell occupancy, completed search cells, spill annotations \\
Route planner & Select a reachable unsearched cell and find a shortest grid path \\
Coverage planner & Generate ordered local waypoints for one cell \\
Motion representation & Convert consecutive waypoints to pivot-and-drive records \\
Execution layer & Currently a GUI animation; physical sensors and motors are absent \\
\bottomrule
\end{tabular}
\end{center}

\section{Map-level route selection}

Let each cell be indexed by $(r,c)$, with $r$ increasing downward and $c$
increasing rightward. The traversable set excludes rack and user-added obstacle
cells. From the robot's current cell $s$, the planner considers every reachable
cell that is neither blocked nor already searched. For each candidate $g$, a
breadth-first search (BFS) finds a shortest path using only the four cardinal
neighbors:
\[
  \mathcal{N}(r,c)=\{(r-1,c),(r+1,c),(r,c-1),(r,c+1)\}.
\]
The selected target is
\[
  g^*=\operatorname*{arg\,min}_{g\in U_{\mathrm{reachable}}}
      \bigl(d_{\mathrm{BFS}}(s,g),\ r_g,\ c_g\bigr),
\]
where the row and column terms make equal-distance choices deterministic.
The route is therefore shortest in grid steps for the selected target; it is
not a globally optimal tour over all remaining cells. Transit cells are
waypoints only. A tile is recorded as searched after its local pattern finishes.

\section{Within-tile coverage}

The local planner uses coordinates $(u,v)$ measured in tile widths from the
tile's upper-left corner. Thus the tile occupies $0\leq u,v\leq1$; points
outside this range are permitted so the robot footprint can extend across tile
boundaries. For the zigzag pattern, the 7-inch lateral dimension gives a
nominal swath width of $7\,\mathrm{in}$. A 2-inch overlap gives lane spacing
\[
  \Delta = 7\,\mathrm{in}-2\,\mathrm{in}=5\,\mathrm{in}.
\]
The generated path has two parallel sweeps on a 12-inch square, joined by
orthogonal repositioning legs. At each corner, the robot pivots in place and
then drives forward along the next segment. The circle-spiral option is a
sampled outward path followed by a retrace toward its start.

The footprint-based swath is only a geometric approximation. A real search
coverage model needs the sensor field of view, desired detection overlap,
localization uncertainty, speed limits, and measured turning behavior.

\section{Global motion-command vector}

The map frame uses feet, with origin at the map's upper-left corner:
\[
  x=c+u,\qquad y=r+v.
\]
The display convention is $+x$ east/right and $+y$ south/down. A segment from
$(x_i,y_i)$ to $(x_{i+1},y_{i+1})$ has length and desired heading
\[
  L_i=\sqrt{(x_{i+1}-x_i)^2+(y_{i+1}-y_i)^2},\qquad
  \theta_i=\operatorname{atan2}(y_{i+1}-y_i,x_{i+1}-x_i).
\]
Given current heading $\theta_{i-1}$, the signed shortest pivot is
\[
  \Delta\theta_i=\operatorname{wrap}_{[-\pi,\pi)}(\theta_i-\theta_{i-1}).
\]
Each entry in the program's \texttt{movement\_command\_vector} contains
\texttt{pivot\_angle\_degrees}, \texttt{global\_heading\_degrees},
\texttt{target\_x\_ft}, \texttt{target\_y\_ft}, and
\texttt{distance\_forward\_ft}. A physical low-level controller would need
to execute the pivot, verify heading, drive toward the target, and verify
position with encoders and/or localization sensors. The current animation
illustrates this sequence but does not provide motor commands.

\section{Spill annotation and saved data}

Spill placement is a two-stage geometric input: choose a map cell, then choose
a point in its enlarged local view. The point is stored as tile index $(r,c)$
and normalized coordinates $(u_s,v_s)\in[0,1]^2$. Its map-frame position is
\[
  (x_s,y_s)=(c+u_s,\ r+v_s)\quad\text{feet}.
\]
Spill annotations are included in the JSON plan alongside rack layout,
completed tile paths, route cells, and the current movement-command vector.
In the simulator, the robot stops and latches when its map position comes
within $3.5\,\mathrm{in}$ of an annotated spill point. The operator must
acknowledge the stop before motion resumes; an acknowledged patch does not
trigger another stop. This proximity rule demonstrates the stop behavior only:
the program has no physical leak sensor, does not prioritize spill tiles, and
does not mark a spill as resolved.

\section{Physical-system development steps}

Before using this plan on hardware, the following interfaces need explicit
design and validation:
\begin{enumerate}
  \item A calibrated floor map and a footprint-aware obstacle-clearance check.
  \item Sensor coverage and leak-detection logic, including confidence and
        false-alarm handling.
  \item A localization source that reports $(x,y,\theta)$ in the map frame.
  \item Closed-loop pivot and forward-drive controllers with speed, stop, and
        emergency-stop limits.
  \item A mission state machine that distinguishes planned, executed,
        interrupted, and verified search segments.
  \item Data logging that records timestamps, robot pose, sensor observations,
        commanded motion, and completion status for each waypoint.
\end{enumerate}

\section{Implementation map}

\begin{itemize}
  \item \texttt{route.py}: rack generation, BFS path reconstruction, and
        nearest reachable unsearched-cell selection.
  \item \texttt{search.py}: zigzag and circle-spiral waypoint generation and
        map search-state storage.
  \item \texttt{main.py}: simulation state, pivot-and-drive vector generation,
        spill-point input, visualization, and JSON export.
\end{itemize}

\section{Software communication and state ownership}

\texttt{main.py} is the coordinator. It imports functions from both planning
modules and passes ordinary Python values between them. The modules do not
communicate through a network, shared database, or hardware bus:
\texttt{route.py} and \texttt{search.py} are called by \texttt{main.py}, and
they do not call each other.

At startup, \texttt{main.py} creates one \texttt{SearchGrid} instance. This
object is the in-memory map-state container. The principal map state is:
\begin{center}
\begin{tabular}{lll}
\toprule
Python field & Type & Meaning \\
\midrule
\texttt{rows}, \texttt{columns} & integers & Map dimensions in cells \\
\texttt{obstacles} & set of $(r,c)$ & Rack cells and user-added blocked cells \\
\texttt{searched} & set of $(r,c)$ & Cells whose complete local pattern finished \\
\texttt{order} & list of $(r,c)$ & Completed cells in completion order \\
\texttt{current\_cell} & $(r,c)$ & Last completed cell used as route origin \\
\texttt{target\_cell} & $(r,c)$ or none & Cell whose local pattern is active \\
\texttt{complete} & Boolean & Whether no reachable work remains \\
\bottomrule
\end{tabular}
\end{center}

These row/column pairs use zero-based Python indices. Thus the upper-left cell
is $(0,0)$ in memory but is reported as row 1, column 1 in saved JSON. A set
stores unique membership efficiently; a list is used for \texttt{order}
because order matters there. The GUI does not own a second map: drawing
functions read these state values and render them on a canvas.

The hand-off for each target proceeds as follows:
\begin{enumerate}
  \item \texttt{main.py} calls
        \texttt{nearest\_unsearched(current, searched, obstacles, rows, columns)}
        in \texttt{route.py}. That function tests eligible cells with
        \texttt{path\_between}, whose BFS returns a list of map cells from
        start through goal.
  \item \texttt{main.py} saves that list as \texttt{planned\_route}; its
        entries after the start become \texttt{travel\_points}. Transit cells
        are visited for travel but remain unsearched.
  \item \texttt{main.py} calls \texttt{make\_search\_path(pattern)} in
        \texttt{search.py}. It receives an ordered list of local $(u,v)$ tile
        fractions. The animation converts a local point to global map feet by
        $x=c+u$ and $y=r+v$ for target cell $(r,c)$.
  \item Tkinter's scheduled \texttt{update\_search} callback advances the
        robot pose a small amount per frame. Once every local point is reached,
        \texttt{complete\_target()} adds the target to \texttt{searched} and
        \texttt{order}; only then is another target selected.
\end{enumerate}

The live robot pose is held separately in \texttt{main.py}: continuous
\texttt{robot\_row}, \texttt{robot\_column}, and \texttt{robot\_heading}
values. This lets the dot move between cell centers. Other active plan state
includes \texttt{path\_points} and \texttt{path\_index} for the current
within-cell pattern, \texttt{travel\_points} for the map route still to
travel, and \texttt{planned\_route} for its complete cell sequence.

\section{Path data structures and persistence}

There are several distinct path representations rather than one universal
path array:
\begin{itemize}
  \item \texttt{planned\_route}: a list of integer $(row,column)$ map cells,
        including the start and target. This is the route between tile centers.
  \item \texttt{path\_points}: a list of floating-point $(u,v)$ fractions for
        the active tile's local coverage pattern. Values may extend outside
        $[0,1]$ to allow the body to pass the tile edge.
  \item \texttt{movement\_command\_vector}: a list of dictionaries generated
        by \texttt{motion\_commands\_from\_pose}. Each describes one upcoming
        pivot-and-forward segment with pivot angle, desired heading, global
        target $(x,y)$ in feet, distance, and action. This is a high-level
        command representation, not motor-controller output.
  \item \texttt{completed\_path\_records}: a list of records appended after a
        tile's local path finishes. Each record identifies the 1-based tile,
        pattern, local waypoints in feet, and segment vectors $(dx,dy)$,
        distance, and heading.
\end{itemize}

The **Save path plan** button calls \texttt{save\_path\_plan}. It serializes
selected in-memory fields as JSON in \texttt{search\_path\_plan.json} beside
\texttt{main.py}. JSON is plain text representing nested dictionaries and
lists, so it can be opened in a text editor and consumed by another program.
The file includes coordinate conventions, robot dimensions, rack layout,
completed paths, the currently planned route and tile path, the current motion
vector, spill patches, and simulated leak events. Spill patches are stored by
integer tile index plus local fractions; their global position follows the
same $x=c+u$, $y=r+v$ conversion.

The save is a snapshot, not automatic persistence: the application does not
currently load the JSON file at startup, and the live tile state (including
\texttt{searched}, \texttt{current\_cell}, and \texttt{target\_cell}) is not
fully serialized as a restorable mission checkpoint. Therefore progress in
memory is lost on exit unless a future load/checkpoint feature is added. Also,
the saved JSON stores completed tile paths and current route/vector data, but
does not by itself prove a physical robot executed them; execution status,
timestamps, and sensor readings would need explicit logging fields.

\end{document}
